%%%%%%%%%%%%%%%%%%%%%%%%%%%%%%%%%%%%%%%%%
% Presentasi Implementasi: Kode Python untuk Quantum Game Theory
% Berdasarkan otak_atik.ipynb
%%%%%%%%%%%%%%%%%%%%%%%%%%%%%%%%%%%%%%%%%

\documentclass[aspectratio=169]{beamer}
\usepackage[utf8]{inputenc}
\usepackage{amsmath,amssymb,physics}
\usepackage{graphicx}
\usepackage{listings}
\usepackage{xcolor}
\usepackage{booktabs}
\usepackage{hyperref}

% Theme
\usetheme{Madrid}
\usecolortheme{seahorse}

% Code listing style
\lstset{
    language=Python,
    basicstyle=\ttfamily\tiny,
    keywordstyle=\color{blue},
    stringstyle=\color{red},
    commentstyle=\color{gray},
    numbers=left,
    numberstyle=\tiny,
    numbersep=5pt,
    breaklines=true,
    showstringspaces=false,
    frame=single,
    backgroundcolor=\color{gray!10}
}

% Judul
\title{Implementasi Quantum Bayesian Game Theory}
\subtitle{Kode Python untuk Analisis Pasar Keuangan}
\author{Presentasi Tugas Akhir}
\date{\today}

\begin{document}

% Slide Judul
\begin{frame}
    \titlepage
\end{frame}

% Outline
\begin{frame}{Outline}
    \tableofcontents
\end{frame}

%%%%%%%%%%%%%%%%%%%%%%%%%%%%%%%%%%%%%%%%%
\section{Import Library dan Data}
%%%%%%%%%%%%%%%%%%%%%%%%%%%%%%%%%%%%%%%%%

\begin{frame}[fragile]{Import Library}
\begin{lstlisting}
import yfinance as yf
import numpy as np
import pandas as pd
import scipy.linalg as la
import networkx as nx
import matplotlib.pyplot as plt
\end{lstlisting}

\vspace{0.3cm}
Library yang digunakan:
\begin{itemize}
    \item \texttt{yfinance}: Mengambil data historis aset
    \item \texttt{numpy}, \texttt{scipy}: Komputasi numerik dan aljabar linear
    \item \texttt{networkx}: Analisis jaringan
    \item \texttt{matplotlib}: Visualisasi
\end{itemize}
\end{frame}

\begin{frame}[fragile]{Fungsi Pengambilan Data}
\begin{lstlisting}
def get_market_data(tickers, period="2y"):
    data = yf.download(tickers, period=period)['Close']
    returns = data.pct_change().dropna()
    return returns

# Eksekusi
tickers = ['GC=F', 'SI=F']  # Gold dan Silver
data_returns = get_market_data(tickers)
\end{lstlisting}

\vspace{0.3cm}
\begin{itemize}
    \item Mengambil data harga penutupan 2 tahun terakhir
    \item Menghitung return harian menggunakan \texttt{pct\_change()}
\end{itemize}
\end{frame}

%%%%%%%%%%%%%%%%%%%%%%%%%%%%%%%%%%%%%%%%%
\section{Konstruksi Matriks Payoff}
%%%%%%%%%%%%%%%%%%%%%%%%%%%%%%%%%%%%%%%%%

\begin{frame}[fragile]{Fungsi Analisis Quantum Game}
\begin{lstlisting}
def analyze_quantum_game(returns, leader_col, follower_col):
    # Tentukan state: 0 jika naik (up), 1 jika turun (down)
    L_state = (returns[leader_col] < 0).astype(int)
    F_state = (returns[follower_col] < 0).astype(int)
    
    # Inisialisasi frekuensi dan akumulasi return
    n_ij = np.zeros((2, 2))
    sum_L = np.zeros((2, 2))
    sum_F = np.zeros((2, 2))
    
    for i in range(len(returns)):
        l, f = L_state.iloc[i], F_state.iloc[i]
        n_ij[l, f] += 1
        sum_L[l, f] += returns[leader_col].iloc[i]
        sum_F[l, f] += returns[follower_col].iloc[i]
\end{lstlisting}
\end{frame}

\begin{frame}[fragile]{Perhitungan Payoff dan Amplitudo}
\begin{lstlisting}
    n_total = n_ij.sum()
    
    # Nilai payoff dengan Laplace Smoothing
    payoff_L = sum_L / (n_ij + 1)
    payoff_F = sum_F / (n_ij + 1)
    
    # Quantum State Coefficients (Amplitudo Probabilitas)
    prob_matrix = n_ij / n_total
    amplitudes = np.sqrt(prob_matrix)  # a00, a01, a10, a11
    
    return payoff_L, payoff_F, amplitudes, prob_matrix
\end{lstlisting}
\end{frame}

\begin{frame}{Hasil Matriks Payoff}
    Output dari analisis Emas (GC=F) vs Perak (SI=F):
    
    \vspace{0.3cm}
    
    \textbf{Matriks Payoff Leader (Gold):}
    \begin{equation*}
        \begin{bmatrix} 0.0104 & 0.0045 \\ -0.0044 & -0.0107 \end{bmatrix}
    \end{equation*}
    
    \textbf{Matriks Payoff Follower (Silver):}
    \begin{equation*}
        \begin{bmatrix} 0.0204 & -0.0085 \\ 0.0071 & -0.0212 \end{bmatrix}
    \end{equation*}
    
    \textbf{Quantum State Coefficients:}
    \begin{equation*}
        \begin{bmatrix} 0.684 & 0.357 \\ 0.334 & 0.541 \end{bmatrix}
    \end{equation*}
\end{frame}

%%%%%%%%%%%%%%%%%%%%%%%%%%%%%%%%%%%%%%%%%
\section{Perhitungan Entropi Kuantum}
%%%%%%%%%%%%%%%%%%%%%%%%%%%%%%%%%%%%%%%%%

\begin{frame}[fragile]{Konstruksi Matriks Densitas}
\begin{lstlisting}
# Matriks Densitas Global
rho_LF = probs / np.sum(probs) 

# Reduced Density Matrix untuk Leader (trace out Follower)
rho_L = np.diag([probs[0,0] + probs[0,1], probs[1,0] + probs[1,1]])

# Reduced Density Matrix untuk Follower (trace out Leader)
rho_F = np.diag([probs[0,0] + probs[1,0], probs[0,1] + probs[1,1]])
\end{lstlisting}
\end{frame}

\begin{frame}[fragile]{Fungsi Von Neumann Entropy}
\begin{lstlisting}
def von_neumann_entropy(rho):
    """Menghitung entropi S = -Tr(rho * log2(rho))"""
    eigvals = la.eigvalsh(rho)
    # Filter nilai nol untuk menghindari error log
    eigvals = eigvals[eigvals > 1e-12]
    return -np.sum(eigvals * np.log2(eigvals))

# Hitung Entropi
s_L = von_neumann_entropy(rho_L)
s_F = von_neumann_entropy(rho_F)
s_LF = von_neumann_entropy(rho_LF)

# Quantum Mutual Information
qmi = s_L + s_F - s_LF
\end{lstlisting}
\end{frame}

\begin{frame}{Hasil Analisis Entropi}
    \begin{table}
        \centering
        \begin{tabular}{lc}
            \toprule
            \textbf{Metrik} & \textbf{Nilai} \\
            \midrule
            Entropi Emas ($S_L$) & 0.9735 \\
            Entropi Perak ($S_F$) & 0.9816 \\
            Quantum Mutual Information ($I$) & 0.9724 \\
            \bottomrule
        \end{tabular}
    \end{table}
    
    \vspace{0.5cm}
    
    \textbf{Interpretasi}: Entanglement sangat kuat. Emas dan Perak memiliki ``jalur informasi'' yang hampir identik.
\end{frame}

%%%%%%%%%%%%%%%%%%%%%%%%%%%%%%%%%%%%%%%%%
\section{Analisis Network Theory}
%%%%%%%%%%%%%%%%%%%%%%%%%%%%%%%%%%%%%%%%%

\begin{frame}[fragile]{Membangun Quantum Financial Network}
\begin{lstlisting}
# Definisikan daftar aset
tickers_list = ['GC=F', 'SI=F', 'CL=F', 'HG=F', 'BTC-USD', 'AAPL']

# Bangun Matriks Adjasensi berdasarkan QMI
G = nx.Graph()

for i, a in enumerate(tickers_list):
    for j, b in enumerate(tickers_list):
        if i < j:
            qmi_val = calculate_qmi_pair(all_data, a, b)
            G.add_edge(a, b, weight=qmi_val)
\end{lstlisting}
\end{frame}

\begin{frame}[fragile]{Eigenvector Centrality Analysis}
\begin{lstlisting}
# Cari Leader Utama dengan Eigenvector Centrality
centrality = nx.eigenvector_centrality(G, weight='weight')
sorted_centrality = sorted(centrality.items(), 
                           key=lambda x: x[1], reverse=True)
\end{lstlisting}

\vspace{0.3cm}

\textbf{Hasil Centrality Analysis (Market Leaders):}
\begin{table}
    \centering
    \begin{tabular}{lc}
        \toprule
        \textbf{Asset} & \textbf{Centrality} \\
        \midrule
        AAPL & 0.4642 \\
        CL=F & 0.4158 \\
        BTC-USD & 0.4054 \\
        HG=F & 0.3980 \\
        GC=F & 0.3812 \\
        SI=F & 0.3789 \\
        \bottomrule
    \end{tabular}
\end{table}
\end{frame}

%%%%%%%%%%%%%%%%%%%%%%%%%%%%%%%%%%%%%%%%%
\section{Konstruksi Hamiltonian}
%%%%%%%%%%%%%%%%%%%%%%%%%%%%%%%%%%%%%%%%%

\begin{frame}[fragile]{Perhitungan Parameter Hamiltonian}
\begin{lstlisting}
# 1. Hitung parameter bias (h) dan interaksi (J)
h_L = (pL[0,0] + pL[0,1]) - (pL[1,0] + pL[1,1])
h_F = (pF[0,0] + pF[1,0]) - (pF[0,1] + pF[1,1])
J_LF = qmi  # Quantum Mutual Information sebagai kekuatan interaksi

# 2. Definisikan Operator Pauli-Z dan Identitas
Z = np.array([[1, 0], [0, -1]])
I_mat = np.eye(2)
\end{lstlisting}
\end{frame}

\begin{frame}[fragile]{Konstruksi Matriks Hamiltonian}
\begin{lstlisting}
# 3. Konstruksi Hamiltonian menggunakan Tensor Product
# H = -hL(Z x I) - hF(I x Z) - J(Z x Z)
H = -h_L * np.kron(Z, I_mat) - h_F * np.kron(I_mat, Z) \
    - J_LF * np.kron(Z, Z)

# 4. Mencari Ground State (Energi Terendah)
eigenvalues, eigenvectors = la.eigh(H)
ground_state_energy = eigenvalues[0]
ground_state_vector = eigenvectors[:, 0]
\end{lstlisting}
\end{frame}

\begin{frame}{Hasil Matriks Hamiltonian}
    \textbf{Matriks Hamiltonian Sistem (4x4):}
    \begin{equation*}
        H = \begin{bmatrix}
            -1.0596 & 0 & 0 & 0 \\
            0 & 0.9997 & 0 & 0 \\
            0 & 0 & 0.9451 & 0 \\
            0 & 0 & 0 & -0.8852
        \end{bmatrix}
    \end{equation*}
    
    \vspace{0.3cm}
    
    \begin{itemize}
        \item \textbf{Energi Ground State}: $-1.0596$
        \item \textbf{Vektor Keadaan Paling Stabil}: $[1, 0, 0, 0]$ (keadaan $\ket{00}$)
    \end{itemize}
\end{frame}

\begin{frame}{Interpretasi Hasil}
    \begin{itemize}
        \item \textbf{Vektor Keadaan Stabil}: $\ket{00}$ (keduanya naik) adalah keadaan paling stabil
        \item \textbf{Energi Sistem}: Jika jauh di atas ground state, pasar dalam kondisi \textbf{instabil}
        \item \textbf{Prediksi Evolusi}: Menggunakan operator $U = e^{-iHt}$ untuk melihat evolusi waktu
    \end{itemize}
    
    \vspace{0.5cm}
    
    Metodologi telah bertransformasi dari analisis data historis menjadi model \textbf{mekanika kuantum stokastik}
\end{frame}

%%%%%%%%%%%%%%%%%%%%%%%%%%%%%%%%%%%%%%%%%
\section{Kesimpulan}
%%%%%%%%%%%%%%%%%%%%%%%%%%%%%%%%%%%%%%%%%

\begin{frame}{Kesimpulan Implementasi}
    Implementasi lengkap mencakup:
    \begin{enumerate}
        \item Pengambilan data pasar dengan \texttt{yfinance}
        \item Konstruksi matriks payoff untuk Game Theory
        \item Perhitungan Von Neumann Entropy
        \item Quantum Mutual Information untuk mengukur entanglement
        \item Network analysis dengan eigenvector centrality
        \item Konstruksi Hamiltonian untuk prediksi evolusi pasar
    \end{enumerate}
\end{frame}

\end{document}
